% Options for packages loaded elsewhere
\PassOptionsToPackage{unicode}{hyperref}
\PassOptionsToPackage{hyphens}{url}
\PassOptionsToPackage{dvipsnames,svgnames*,x11names*}{xcolor}
%
\documentclass[
  UTF8,fontset=none,zihao=5]{ctexart}
\usepackage{amsmath,amssymb}
\usepackage{lmodern}
\usepackage{iftex}
\ifPDFTeX
  \usepackage[T1]{fontenc}
  \usepackage[utf8]{inputenc}
  \usepackage{textcomp} % provide euro and other symbols
\else % if luatex or xetex
  \usepackage{unicode-math}
  \defaultfontfeatures{Scale=MatchLowercase}
  \defaultfontfeatures[\rmfamily]{Ligatures=TeX,Scale=1}
\fi
% Use upquote if available, for straight quotes in verbatim environments
\IfFileExists{upquote.sty}{\usepackage{upquote}}{}
\IfFileExists{microtype.sty}{% use microtype if available
  \usepackage[]{microtype}
  \UseMicrotypeSet[protrusion]{basicmath} % disable protrusion for tt fonts
}{}
\makeatletter
\@ifundefined{KOMAClassName}{% if non-KOMA class
  \IfFileExists{parskip.sty}{%
    \usepackage{parskip}
  }{% else
    \setlength{\parindent}{0pt}
    \setlength{\parskip}{6pt plus 2pt minus 1pt}}
}{% if KOMA class
  \KOMAoptions{parskip=half}}
\makeatother
\usepackage{xcolor}
\IfFileExists{xurl.sty}{\usepackage{xurl}}{} % add URL line breaks if available
\IfFileExists{bookmark.sty}{\usepackage{bookmark}}{\usepackage{hyperref}}
\hypersetup{
  pdftitle={算法设计与分析 Day4：随机快速排序、排序下界与三道作业详解},
  pdfauthor={RedMoon},
  colorlinks=true,
  linkcolor={NoteGreen},
  filecolor={Maroon},
  citecolor={Blue},
  urlcolor={NoteGreen},
  pdfcreator={LaTeX via pandoc}}
\urlstyle{same} % disable monospaced font for URLs
\usepackage{color}
\usepackage{fancyvrb}
\newcommand{\VerbBar}{|}
\newcommand{\VERB}{\Verb[commandchars=\\\{\}]}
\DefineVerbatimEnvironment{Highlighting}{Verbatim}{commandchars=\\\{\}}
% Add ',fontsize=\small' for more characters per line
\newenvironment{Shaded}{}{}
\newcommand{\AlertTok}[1]{\textcolor[rgb]{1.00,0.00,0.00}{\textbf{#1}}}
\newcommand{\AnnotationTok}[1]{\textcolor[rgb]{0.38,0.63,0.69}{\textbf{\textit{#1}}}}
\newcommand{\AttributeTok}[1]{\textcolor[rgb]{0.49,0.56,0.16}{#1}}
\newcommand{\BaseNTok}[1]{\textcolor[rgb]{0.25,0.63,0.44}{#1}}
\newcommand{\BuiltInTok}[1]{#1}
\newcommand{\CharTok}[1]{\textcolor[rgb]{0.25,0.44,0.63}{#1}}
\newcommand{\CommentTok}[1]{\textcolor[rgb]{0.38,0.63,0.69}{\textit{#1}}}
\newcommand{\CommentVarTok}[1]{\textcolor[rgb]{0.38,0.63,0.69}{\textbf{\textit{#1}}}}
\newcommand{\ConstantTok}[1]{\textcolor[rgb]{0.53,0.00,0.00}{#1}}
\newcommand{\ControlFlowTok}[1]{\textcolor[rgb]{0.00,0.44,0.13}{\textbf{#1}}}
\newcommand{\DataTypeTok}[1]{\textcolor[rgb]{0.56,0.13,0.00}{#1}}
\newcommand{\DecValTok}[1]{\textcolor[rgb]{0.25,0.63,0.44}{#1}}
\newcommand{\DocumentationTok}[1]{\textcolor[rgb]{0.73,0.13,0.13}{\textit{#1}}}
\newcommand{\ErrorTok}[1]{\textcolor[rgb]{1.00,0.00,0.00}{\textbf{#1}}}
\newcommand{\ExtensionTok}[1]{#1}
\newcommand{\FloatTok}[1]{\textcolor[rgb]{0.25,0.63,0.44}{#1}}
\newcommand{\FunctionTok}[1]{\textcolor[rgb]{0.02,0.16,0.49}{#1}}
\newcommand{\ImportTok}[1]{#1}
\newcommand{\InformationTok}[1]{\textcolor[rgb]{0.38,0.63,0.69}{\textbf{\textit{#1}}}}
\newcommand{\KeywordTok}[1]{\textcolor[rgb]{0.00,0.44,0.13}{\textbf{#1}}}
\newcommand{\NormalTok}[1]{#1}
\newcommand{\OperatorTok}[1]{\textcolor[rgb]{0.40,0.40,0.40}{#1}}
\newcommand{\OtherTok}[1]{\textcolor[rgb]{0.00,0.44,0.13}{#1}}
\newcommand{\PreprocessorTok}[1]{\textcolor[rgb]{0.74,0.48,0.00}{#1}}
\newcommand{\RegionMarkerTok}[1]{#1}
\newcommand{\SpecialCharTok}[1]{\textcolor[rgb]{0.25,0.44,0.63}{#1}}
\newcommand{\SpecialStringTok}[1]{\textcolor[rgb]{0.73,0.40,0.53}{#1}}
\newcommand{\StringTok}[1]{\textcolor[rgb]{0.25,0.44,0.63}{#1}}
\newcommand{\VariableTok}[1]{\textcolor[rgb]{0.10,0.09,0.49}{#1}}
\newcommand{\VerbatimStringTok}[1]{\textcolor[rgb]{0.25,0.44,0.63}{#1}}
\newcommand{\WarningTok}[1]{\textcolor[rgb]{0.38,0.63,0.69}{\textbf{\textit{#1}}}}
\usepackage{longtable,booktabs,array}
\usepackage{calc} % for calculating minipage widths
% Correct order of tables after \paragraph or \subparagraph
\usepackage{etoolbox}
\makeatletter
\patchcmd\longtable{\par}{\if@noskipsec\mbox{}\fi\par}{}{}
\makeatother
% Allow footnotes in longtable head/foot
\IfFileExists{footnotehyper.sty}{\usepackage{footnotehyper}}{\usepackage{footnote}}
\makesavenoteenv{longtable}
\usepackage{graphicx}
\makeatletter
\def\maxwidth{\ifdim\Gin@nat@width>\linewidth\linewidth\else\Gin@nat@width\fi}
\def\maxheight{\ifdim\Gin@nat@height>\textheight\textheight\else\Gin@nat@height\fi}
\makeatother
% Scale images if necessary, so that they will not overflow the page
% margins by default, and it is still possible to overwrite the defaults
% using explicit options in \includegraphics[width, height, ...]{}
\setkeys{Gin}{width=\maxwidth,height=\maxheight,keepaspectratio}
% Set default figure placement to htbp
\makeatletter
\def\fps@figure{htbp}
\makeatother
\setlength{\emergencystretch}{3em} % prevent overfull lines
\providecommand{\tightlist}{%
  \setlength{\itemsep}{0pt}\setlength{\parskip}{0pt}}
\setcounter{secnumdepth}{-\maxdimen} % remove section numbering
\usepackage[a4paper,top=20mm,bottom=20mm,left=20mm,right=20mm,headheight=14pt,headsep=7mm]{geometry}
\usepackage{fancyhdr}
\usepackage{needspace}
\usepackage{multicol}
\usepackage{microtype}
\usepackage{fvextra}
\usepackage{caption}
\usepackage{framed}
\definecolor{NoteGreen}{HTML}{17685F}
\definecolor{NoteRed}{HTML}{A52B34}
\definecolor{NoteGray}{HTML}{667078}
\definecolor{shadecolor}{HTML}{F3F6F6}
\renewenvironment{Shaded}{\begin{snugshade}}{\end{snugshade}}
\hypersetup{pdflang=zh-CN}
\setmainfont{TeX Gyre Pagella}
\setsansfont{TeX Gyre Heros}
\setmonofont{Consolas}[Scale=0.89]
\setCJKmainfont{SimSun}[AutoFakeBold=2]
\setCJKsansfont{Microsoft YaHei}
\setCJKmonofont{Microsoft YaHei}
\linespread{1.18}
\setlength{\parindent}{0pt}
\setlength{\parskip}{5pt plus 1pt minus 1pt}
\setlength{\emergencystretch}{2em}
\setlength{\columnsep}{8mm}
\setlength{\LTpre}{5pt}
\setlength{\LTpost}{7pt}
\renewcommand{\arraystretch}{1.18}
\captionsetup{font=small,labelfont=bf,skip=7pt}
\ctexset{
  section={format=\Large\sffamily\bfseries\color{NoteGreen},beforeskip=17pt,afterskip=9pt},
  subsection={format=\large\sffamily\bfseries,beforeskip=12pt,afterskip=6pt},
  subsubsection={format=\normalsize\sffamily\bfseries,beforeskip=9pt,afterskip=5pt}
}
\let\OldSection\section
\renewcommand{\section}{\Needspace{5\baselineskip}\OldSection}
\let\OldSubsection\subsection
\renewcommand{\subsection}{\Needspace{4\baselineskip}\OldSubsection}
\pagestyle{fancy}
\fancyhf{}
\fancyhead[L]{\small\sffamily\color{NoteGray}算法设计与分析 \quad Day4}
\fancyhead[R]{\small\sffamily\color{NoteGray}随机排序与比较下界}
\fancyfoot[L]{\footnotesize\sffamily\color{NoteGray}RedMoon \quad 2026.10.15}
\fancyfoot[R]{\small\thepage}
\renewcommand{\headrulewidth}{0.3pt}
\renewcommand{\footrulewidth}{0pt}
\setcounter{tocdepth}{1}
\setcounter{secnumdepth}{-1}
\makeatletter
\renewcommand{\maketitle}{
  \thispagestyle{fancy}
  {\small\sffamily\color{NoteRed}COURSE NOTES / DAY 04\par}
  \vspace{5pt}
  {\fontsize{23}{30}\selectfont\sffamily\bfseries 算法设计与分析\quad Day4\par}
  \vspace{4pt}
  {\large\sffamily 随机快速排序、排序下界与三道作业详解\par}
  \vspace{7pt}
  {\small\color{NoteGray}RedMoon \quad 排序进阶 \quad 2026 年 10 月 15 日\par}
  \vspace{5pt}
  {\color{NoteGreen}\hrule height 0.8pt}
  \vspace{7pt}
}
\let\OldTableOfContents\tableofcontents
\renewcommand{\tableofcontents}{
  \begingroup
  \small
  \begin{multicols}{2}
  \OldTableOfContents
  \end{multicols}
  \endgroup
  \vspace{4pt}
}
\AtBeginEnvironment{longtable}{\small}
\AtBeginEnvironment{Shaded}{\small}
\fvset{breaklines=true,breakanywhere=true,fontsize=\small}
\makeatother
\ifLuaTeX
  \usepackage{selnolig}  % disable illegal ligatures
\fi

\title{算法设计与分析 Day4：随机快速排序、排序下界与三道作业详解}
\author{RedMoon}
\date{2026-10-15}

\begin{document}
\maketitle

{
\hypersetup{linkcolor=NoteGreen}
\setcounter{tocdepth}{1}
\tableofcontents
}
\begin{quote}
RedMoon · 算法设计与分析 · Day4
\end{quote}

快速排序的一次划分会把枢轴放到最终位置。剩下的工作量，取决于枢轴两侧各有多少元素。本次课从划分规模出发，推导随机快速排序的期望时间，再结合堆排序、计数排序和决策树，讨论排序算法利用的信息及其复杂度下界。

文末详解三道作业：线性时间最大子数组、用最小堆合并有序表，以及快速排序与插入排序的组合。每题包含题意、算法、执行示例、正确性证明和复杂度分析。

记 \(\lg n=\log_2 n\)，\(\ln n\) 为自然对数；区间 \(A[l..r]\)
包含两端。Python 代码使用 \textbf{0-based
下标}，返回的最大子数组区间也包含左右端点。时间分析默认一次比较、数组访问及机器字范围内的算术运算花费常数时间；涉及输入编码长度时，会单独说明。

\hypertarget{ux4e00ux8bfeux5802ux70edux8eabux591aux6570ux5143ux7d20ux4e3aux4ec0ux4e48ux53efux4ee5ux4e24ux4e24ux62b5ux6d88}{%
\section{一、课堂热身：多数元素为什么可以两两抵消}\label{ux4e00ux8bfeux5802ux70edux8eabux591aux6570ux5143ux7d20ux4e3aux4ec0ux4e48ux53efux4ee5ux4e24ux4e24ux62b5ux6d88}}

\hypertarget{ux591aux6570ux5143ux7d20ux4e0eux4f17ux6570ux7684ux533aux522b}{%
\subsection{1.
多数元素与众数的区别}\label{ux591aux6570ux5143ux7d20ux4e0eux4f17ux6570ux7684ux533aux522b}}

课件第 2 页的题目是 LeetCode
169「多数元素」：给定一个非空数组，保证存在出现次数 \textbf{严格大于}
\(\lfloor n/2\rfloor\) 的元素，找出它。

众数只要求出现次数最多，多数元素则要求占据一半以上。例如 \([1,1,2,2,3]\)
中，1 和 2
都是众数，但没有多数元素。严格多数至多只有一个，因为两个数的出现次数不可能都超过
\(n/2\)。

直接统计频次即可求解。用哈希表记录，期望时间为 \(O(n)\)，最坏空间为
\(O(n)\)；若已知整数值域为 \(1..K\)，也可以使用长度为 \(K\)
的计数数组，时间为 \(O(n+K)\)，空间为
\(O(K)\)。值域越大，初始化计数数组的成本越高。

\hypertarget{boyermoore-ux6295ux7968ux6cd5}{%
\subsection{2. Boyer--Moore
投票法}\label{boyermoore-ux6295ux7968ux6cd5}}

从数组中删去两个值不同的元素，原来的多数元素仍然是剩余数组中的多数元素。

设多数元素为 \(x\)，原来出现 \(c>n/2\) 次。若删去的一对中包含
\(x\)，它还剩 \(c-1\) 次，且 \(c-1>(n-2)/2\)；若这一对都不是
\(x\)，它的次数不变，占比反而更高。因此，不断抵消后留下的候选值就是
\(x\)。

实现时维护一个候选值 \texttt{candidate} 和未被抵消的数量
\texttt{count}：

\Needspace{106pt}

\begin{Shaded}
\begin{Highlighting}[]
\KeywordTok{def}\NormalTok{ majority\_element(a):}
    \ControlFlowTok{if} \KeywordTok{not}\NormalTok{ a:}
        \ControlFlowTok{raise} \PreprocessorTok{ValueError}\NormalTok{(}\StringTok{"input must be nonempty"}\NormalTok{)}
\NormalTok{    candidate, count }\OperatorTok{=} \VariableTok{None}\NormalTok{, }\DecValTok{0}
    \ControlFlowTok{for}\NormalTok{ x }\KeywordTok{in}\NormalTok{ a:}
        \ControlFlowTok{if}\NormalTok{ count }\OperatorTok{==} \DecValTok{0}\NormalTok{:}
\NormalTok{            candidate }\OperatorTok{=}\NormalTok{ x}
\NormalTok{        count }\OperatorTok{+=} \DecValTok{1} \ControlFlowTok{if}\NormalTok{ x }\OperatorTok{==}\NormalTok{ candidate }\ControlFlowTok{else} \OperatorTok{{-}}\DecValTok{1}
    \ControlFlowTok{return}\NormalTok{ candidate}
\end{Highlighting}
\end{Shaded}

以 \([2,2,1,1,1,2,2]\) 为例：

\begin{longtable}[]{@{}
  >{\raggedright\arraybackslash}p{(\columnwidth - 8\tabcolsep) * \real{0.10}}
  >{\raggedleft\arraybackslash}p{(\columnwidth - 8\tabcolsep) * \real{0.13}}
  >{\raggedleft\arraybackslash}p{(\columnwidth - 8\tabcolsep) * \real{0.21}}
  >{\raggedleft\arraybackslash}p{(\columnwidth - 8\tabcolsep) * \real{0.13}}
  >{\raggedright\arraybackslash}p{(\columnwidth - 8\tabcolsep) * \real{0.43}}@{}}
\toprule
下标 & 读入值 & 处理后的候选值 & \texttt{count} & 发生的操作 \\
\midrule
\endhead
0 & 2 & 2 & 1 & 空状态建立候选值 \\
1 & 2 & 2 & 2 & 同值累积 \\
2 & 1 & 2 & 1 & 抵消一对不同值 \\
3 & 1 & 2 & 0 & 当前候选全部抵消 \\
4 & 1 & 1 & 1 & 重新建立候选值 \\
5 & 2 & 1 & 0 & 再次抵消 \\
6 & 2 & 2 & 1 & 剩下候选值 2 \\
\bottomrule
\end{longtable}

\texttt{count}
记录尚未被抵消的候选数量，不代表候选值在整个数组中的出现次数。算法只扫描一次，时间
\(\Theta(n)\)，辅助空间 \(O(1)\)。

若题目不保证多数元素存在，还需在第二遍扫描中统计候选值，检查它是否出现超过
\(\lfloor n/2\rfloor\) 次。例如 \([1,2,3]\)
第一遍也会产生候选值，但数组中没有多数元素。

\hypertarget{ux4e8cux5206ux6cbbux6392ux5e8fux5f52ux5e76ux6392ux5e8fux4e0eux5febux901fux6392ux5e8fux7684ux5206ux5de5}{%
\section{二、分治排序：归并排序与快速排序的分工}\label{ux4e8cux5206ux6cbbux6392ux5e8fux5f52ux5e76ux6392ux5e8fux4e0eux5febux901fux6392ux5e8fux7684ux5206ux5de5}}

排序的输出既要按非降序排列，也要与输入包含完全相同的元素及重复次数。正确性证明需要同时覆盖有序性与元素保留，避免遗漏或增加元素。

课堂先区分了问题、算法和分析。排序规定要解决什么问题，归并排序与快速排序给出具体做法，正确性和运行时间分析则回答这些做法是否正确、需要多少代价。同一个问题可以有多种算法，分析时也可以采用不同模型。

\begin{longtable}[]{@{}
  >{\raggedright\arraybackslash}p{(\columnwidth - 4\tabcolsep) * \real{0.16}}
  >{\raggedright\arraybackslash}p{(\columnwidth - 4\tabcolsep) * \real{0.42}}
  >{\raggedright\arraybackslash}p{(\columnwidth - 4\tabcolsep) * \real{0.42}}@{}}
\toprule
分治环节 & 归并排序 & 快速排序 \\
\midrule
\endhead
边界 & 长度不超过 1 时返回 & 长度不超过 1 时返回 \\
分 & 按下标切成两半，切分本身简单 & 扫描元素，按枢轴值划分 \\
治 & 排序左右两半 & 排序枢轴两侧 \\
合 & 线性合并两个有序区间 & 枢轴已经就位，无须再线性合并 \\
子问题规模 & 两侧约为 \(n/2\) & 可能从 \(0,n-1\) 到接近均分 \\
\bottomrule
\end{longtable}

归并排序的主要线性工作在合并阶段，快速排序则在划分阶段。分析快速排序时，需要继续考察划分后两个子问题的规模。

\hypertarget{ux4e09ux968fux673aux5febux901fux6392ux5e8fux5212ux5206ux4e0eux5b9eux73b0}{%
\section{三、随机快速排序：划分与实现}\label{ux4e09ux968fux673aux5febux901fux6392ux5e8fux5212ux5206ux4e0eux5b9eux73b0}}

\hypertarget{ux67a2ux8f74ux5fc5ux987bux9000ux51faux540eux7eedux9012ux5f52}{%
\subsection{1.
枢轴必须退出后续递归}\label{ux67a2ux8f74ux5fc5ux987bux9000ux51faux540eux7eedux9012ux5f52}}

先从数组中均匀随机选取一个位置，以该位置的元素为枢轴
\(x\)。单独取出这次出现，再将其余元素分成：

\[
B=\{a:a\le x\},\qquad C=\{a:a>x\}.
\]

递归结果为：

\[
\operatorname{sort}(B)+[x]+\operatorname{sort}(C).
\]

这里的加号表示顺序拼接。枢轴本身不再进入 \(B\) 或
\(C\)，所以两侧规模之和是 \(n-1\)，每个递归子问题都严格变小。

数组含重复值时，也必须排除选作枢轴的这次出现。否则，把枢轴放回「小于等于」的一侧后，全相等数组可能原样进入下一次递归，无法结束。原地划分会固定一个枢轴位置，使后续区间严格缩小。

\hypertarget{ux539fux5730ux5212ux5206ux7684ux56dbux4e2aux533aux57df}{%
\subsection{2.
原地划分的四个区域}\label{ux539fux5730ux5212ux5206ux7684ux56dbux4e2aux533aux57df}}

课件第 8 页使用 Lomuto 划分。先把随机选中的枢轴换到 \(A[r]\)，令
\(x=A[r]\)，再扫描 \(A[p..r-1]\)。

循环每次开始时维护：

\begin{longtable}[]{@{}
  >{\raggedright\arraybackslash}p{(\columnwidth - 2\tabcolsep) * \real{0.27}}
  >{\raggedright\arraybackslash}p{(\columnwidth - 2\tabcolsep) * \real{0.73}}@{}}
\toprule
区域 & 语义 \\
\midrule
\endhead
\(A[p..i]\) & 已处理且全部 \(\le x\) \\
\(A[i+1..j-1]\) & 已处理且全部 \(>x\) \\
\(A[j..r-1]\) & 尚未检查 \\
\(A[r]\) & 暂存枢轴 \(x\) \\
\bottomrule
\end{longtable}

若 \(A[j]\le x\)，先令 \(i\) 加 1，再交换 \(A[i]\) 与
\(A[j]\)，把这个元素纳入左区间；若 \(A[j]>x\)，则直接推进扫描位置。

扫描结束后，交换 \(A[i+1]\) 与 \(A[r]\)。令 \(q=i+1\)，得到：

\[
A[p..q-1]\le A[q],\qquad A[q+1..r]>A[q].
\]

\Needspace{156pt}

\begin{Shaded}
\begin{Highlighting}[]
\ImportTok{from}\NormalTok{ random }\ImportTok{import}\NormalTok{ Random}


\KeywordTok{def}\NormalTok{ randomized\_partition(a, p, r, rng):}
\NormalTok{    pivot\_index }\OperatorTok{=}\NormalTok{ rng.randrange(p, r }\OperatorTok{+} \DecValTok{1}\NormalTok{)}
\NormalTok{    a[pivot\_index], a[r] }\OperatorTok{=}\NormalTok{ a[r], a[pivot\_index]}
\NormalTok{    x }\OperatorTok{=}\NormalTok{ a[r]}
\NormalTok{    i }\OperatorTok{=}\NormalTok{ p }\OperatorTok{{-}} \DecValTok{1}
    \ControlFlowTok{for}\NormalTok{ j }\KeywordTok{in} \BuiltInTok{range}\NormalTok{(p, r):}
        \ControlFlowTok{if}\NormalTok{ a[j] }\OperatorTok{\textless{}=}\NormalTok{ x:}
\NormalTok{            i }\OperatorTok{+=} \DecValTok{1}
\NormalTok{            a[i], a[j] }\OperatorTok{=}\NormalTok{ a[j], a[i]}
\NormalTok{    a[i }\OperatorTok{+} \DecValTok{1}\NormalTok{], a[r] }\OperatorTok{=}\NormalTok{ a[r], a[i }\OperatorTok{+} \DecValTok{1}\NormalTok{]}
    \ControlFlowTok{return}\NormalTok{ i }\OperatorTok{+} \DecValTok{1}
\end{Highlighting}
\end{Shaded}

\Needspace{146pt}

\begin{Shaded}
\begin{Highlighting}[]
\KeywordTok{def}\NormalTok{ quicksort(a, rng}\OperatorTok{=}\VariableTok{None}\NormalTok{):}
    \ControlFlowTok{if}\NormalTok{ rng }\KeywordTok{is} \VariableTok{None}\NormalTok{:}
\NormalTok{        rng }\OperatorTok{=}\NormalTok{ Random()}

    \KeywordTok{def}\NormalTok{ sort(p, r):}
        \ControlFlowTok{if}\NormalTok{ p }\OperatorTok{\textgreater{}=}\NormalTok{ r:}
            \ControlFlowTok{return}
\NormalTok{        q }\OperatorTok{=}\NormalTok{ randomized\_partition(a, p, r, rng)}
\NormalTok{        sort(p, q }\OperatorTok{{-}} \DecValTok{1}\NormalTok{)}
\NormalTok{        sort(q }\OperatorTok{+} \DecValTok{1}\NormalTok{, r)}

\NormalTok{    sort(}\DecValTok{0}\NormalTok{, }\BuiltInTok{len}\NormalTok{(a) }\OperatorTok{{-}} \DecValTok{1}\NormalTok{)}
    \ControlFlowTok{return}\NormalTok{ a}
\end{Highlighting}
\end{Shaded}

\hypertarget{ux8bfeux4ef6ux793aux4f8bux7684ux4ea4ux6362ux8fc7ux7a0b}{%
\subsection{3.
课件示例的交换过程}\label{ux8bfeux4ef6ux793aux4f8bux7684ux4ea4ux6362ux8fc7ux7a0b}}

输入 \([40,30,10,20]\)，取末尾的 20 为枢轴。使用 0-based 下标，初始
\(i=-1\)：

\begin{longtable}[]{@{}
  >{\raggedright\arraybackslash}p{(\columnwidth - 4\tabcolsep) * \real{0.15}}
  >{\raggedright\arraybackslash}p{(\columnwidth - 4\tabcolsep) * \real{0.47}}
  >{\raggedright\arraybackslash}p{(\columnwidth - 4\tabcolsep) * \real{0.38}}@{}}
\toprule
扫描阶段 & 比较结果与动作 & 数组 \\
\midrule
\endhead
\(j=0\) & \(40>20\)，不交换 & \([40,30,10,20]\) \\
\(j=1\) & \(30>20\)，不交换 & \([40,30,10,20]\) \\
\(j=2\) & \(10\le20\)，令 \(i=0\)，交换位置 0 和 2 &
\([10,30,40,20]\) \\
扫描结束 & 交换位置 \(i+1=1\) 和末尾枢轴 & \([10,20,40,30]\) \\
\bottomrule
\end{longtable}

返回 \(q=1\)。20 已经就位，但右侧 \([40,30]\)
还没有排序。一次划分只保证枢轴两侧满足大小关系。

\hypertarget{ux6b63ux786eux6027ux4e0eux7a7aux95f4}{%
\subsection{4.
正确性与空间}\label{ux6b63ux786eux6027ux4e0eux7a7aux95f4}}

用循环不变式证明划分正确。初始化时，两个已处理区域均为空；扫描中，每次交换将符合条件的元素放进左区间，并保持右区间的性质；扫描结束时，未知区域为空，最后一次交换将枢轴放到两区之间。

再对区间长度归纳，即可证明快速排序正确。长度不超过 1
时，区间已有序。更长的区间经过划分后，枢轴两侧满足大小关系；由归纳假设，两个更短区间递归排序后也各自有序，拼接起来便得到整个有序区间。全过程只交换已有元素，因此输入的多重集合保持不变。

长度为 \(m=r-p+1\) 的划分需要 \(m-1\) 次关键字比较，时间
\(\Theta(m)\)，辅助空间
\(O(1)\)。上面的快速排序仍有递归栈：互异元素、随机枢轴时，期望栈深为
\(O(\log n)\)，最坏可达
\(O(n)\)。若每次只递归较短一侧、用循环处理较长一侧，可以把栈空间最坏限制在
\(O(\log n)\)。

\hypertarget{ux91cdux590dux503cux4e0eux7a33ux5b9aux6027}{%
\subsection{5.
重复值与稳定性}\label{ux91cdux590dux503cux4e0eux7a33ux5b9aux6027}}

Lomuto 使用 \texttt{\textless{}=} 判断。所有元素相等时，每次划分都把其余
\(m-1\) 个元素留在左侧，无论随机选中哪个位置，时间都是
\(\Theta(n^2)\)。随机化改变了枢轴的选择方式，但这个划分规则仍会把所有相等元素放到同一侧。

课件之外，还可以用三路划分处理重复值：将元素分成 \(<x\)、\(=x\)、\(>x\)
三区，只递归两侧；所有值相等时，一次线性扫描即可结束。后续期望递推假设
\textbf{元素互异}，此时每个枢轴秩恰好对应一种左右规模。

这种交换式快速排序通常不稳定。例如给相同键附上标签，\([(2,a),(2,b),(1,c)]\)
以 1 划分后成为 \([(1,c),(2,b),(2,a)]\)，两个键为 2 的记录已交换先后。

\hypertarget{ux56dbux8f93ux5165ux968fux673aux9009ux62e9ux4e0eux8fd0ux884cux65f6ux95f4}{%
\section{四、输入、随机选择与运行时间}\label{ux56dbux8f93ux5165ux968fux673aux9009ux62e9ux4e0eux8fd0ux884cux65f6ux95f4}}

\hypertarget{ux8fd0ux884cux65f6ux95f4ux6709ux4e24ux4e2aux81eaux53d8ux91cf}{%
\subsection{1.
运行时间有两个自变量}\label{ux8fd0ux884cux65f6ux95f4ux6709ux4e24ux4e2aux81eaux53d8ux91cf}}

记 \(T(A,R)\) 为输入数组 \(A\)、随机选择序列 \(R\)
下的运行时间。即使固定 \(A\)，不同的枢轴序列也会产生不同的递归树。

「随机快速排序的期望时间」通常对算法的随机选择取期望：

\[
\max_{|A|=n}\mathbb E_R[T(A,R)].
\]

若元素互异，这个期望为
\(O(n\log n)\)。分析时先固定任意输入，再对枢轴的随机选择取平均，无须假设输入本身是随机排列。

平均情况 \(\mathbb E_A[T(A)]\)
则对输入取平均，需要先给定输入分布。它可能与随机算法的期望时间具有相同的渐近阶，但两者取平均的对象不同。

\hypertarget{ux4e0dux5e26ux6709ux5e8fux6027ux9884ux68c0ux67e5ux7684ux666eux901aux7248ux672c}{%
\subsection{2.
不带有序性预检查的普通版本}\label{ux4e0dux5e26ux6709ux5e8fux6027ux9884ux68c0ux67e5ux7684ux666eux901aux7248ux672c}}

若每次都接近均分：

\[
T(n)=2T(n/2)+\Theta(n)=\Theta(n\log n).
\]

若每次都选到当前最小值或最大值：

\[
T(n)=T(n-1)+\Theta(n)=\Theta(n^2).
\]

后一种情况下，扫描规模累加为
\(n+(n-1)+\cdots+2\)。随机快速排序也可能连续选到这样的枢轴，因此最坏时间仍为二次量级，未降至
\(O(n\log n)\)。

\hypertarget{ux52a0ux5165ux6709ux5e8fux6027ux68c0ux67e5ux540eux7684ux7248ux672c}{%
\subsection{3.
加入有序性检查后的版本}\label{ux52a0ux5165ux6709ux5e8fux6027ux68c0ux67e5ux540eux7684ux7248ux672c}}

课件第 11
页另加了一步：先扫描数组，发现已经有序就返回。对这个变体，已排序输入的运行时间是
\(\Theta(n)\)，因为只需完成一次检查。

带预检查版本在有序输入上花费的时间，与普通快速排序在均衡划分下的
\(\Theta(n\log n)\) 对应不同条件。几种情况如下：

\begin{longtable}[]{@{}
  >{\raggedright\arraybackslash}p{(\columnwidth - 2\tabcolsep) * \real{0.70}}
  >{\raggedright\arraybackslash}p{(\columnwidth - 2\tabcolsep) * \real{0.30}}@{}}
\toprule
所讨论的情况 & 时间 \\
\midrule
\endhead
普通版本，划分始终接近均衡 & \(\Theta(n\log n)\) \\
普通版本，划分始终极端不均衡 & \(\Theta(n^2)\) \\
带预检查版本，输入原本有序 & \(\Theta(n)\) \\
互异元素的随机版本，对随机枢轴取期望 & \(\Theta(n\log n)\) \\
\bottomrule
\end{longtable}

预检查本身需要线性时间。下一节沿用课件的规范化递推，分析不提前返回的普通随机版本。

\hypertarget{ux4e94ux968fux673aux5febux901fux6392ux5e8fux7684ux671fux671bux65f6ux95f4ux63a8ux5bfc}{%
\section{五、随机快速排序的期望时间推导}\label{ux4e94ux968fux673aux5febux901fux6392ux5e8fux7684ux671fux671bux65f6ux95f4ux63a8ux5bfc}}

\hypertarget{ux6309ux67a2ux8f74ux79e9ux5206ux7c7b}{%
\subsection{1.
按枢轴秩分类}\label{ux6309ux67a2ux8f74ux79e9ux5206ux7c7b}}

设 \(T(n)\) 为规模 \(n\)
的期望代价。元素互异、从当前子数组均匀选择枢轴时，枢轴可能有
\(0,1,\ldots,n-1\) 个更小的元素，每种情况概率都是 \(1/n\)。

当左侧规模为 \(i\) 时，右侧规模为
\(n-1-i\)。将一次线性处理的代价规范化写成 \(n\)，得到：

\[
T(n)=\frac1n\sum_{i=0}^{n-1}
\bigl(T(i)+T(n-1-i)\bigr)+n.
\]

两部分求和包含相同的项，只是次序相反，因此：

\[
T(n)=\frac2n\sum_{i=0}^{n-1}T(i)+n.
\]

若保留具体实现中的常数项和线性项系数，推导方式与渐近阶不变。边界代价
\(T(0),T(1)\)
均为常数；相邻规模作差时，要从两个递推式都成立的规模开始，再把有限个初值并入常数。

\hypertarget{ux4e58ux4ee5ux89c4ux6a21ux518dux5bf9ux76f8ux90bbux5f0fux5b50ux4f5cux5dee}{%
\subsection{2.
乘以规模，再对相邻式子作差}\label{ux4e58ux4ee5ux89c4ux6a21ux518dux5bf9ux76f8ux90bbux5f0fux5b50ux4f5cux5dee}}

两边乘以 \(n\)：

\[
nT(n)=2\sum_{i=0}^{n-1}T(i)+n^2.
\]

把规模换为 \(n-1\)：

\[
(n-1)T(n-1)=2\sum_{i=0}^{n-2}T(i)+(n-1)^2.
\]

两式相减，求和只留下最后一项：

\[
nT(n)-(n-1)T(n-1)=2T(n-1)+2n-1,
\]

于是：

\[
nT(n)=(n+1)T(n-1)+2n-1.
\]

求和项由此消去，递推只剩相邻两个规模。

\hypertarget{ux9664ux4ee5-nn1ux5f97ux5230ux53efux7d2fux52a0ux7684ux5f62ux5f0f}{%
\subsection{\texorpdfstring{3. 除以
\(n(n+1)\)，得到可累加的形式}{3. 除以 n(n+1)，得到可累加的形式}}\label{ux9664ux4ee5-nn1ux5f97ux5230ux53efux7d2fux52a0ux7684ux5f62ux5f0f}}

\[
\frac{T(n)}{n+1}
=\frac{T(n-1)}n+\frac{2n-1}{n(n+1)}
\le\frac{T(n-1)}n+\frac2{n+1}.
\]

令 \(S(n)=T(n)/(n+1)\)，就有
\(S(n)\le S(n-1)+2/(n+1)\)。展开至一个固定初值 \(n_0\)：

\[
\frac{T(n)}{n+1}
\le \frac{T(n_0)}{n_0+1}
+2\sum_{j=n_0+1}^{n}\frac1{j+1}
=O(\log n).
\]

因此：

\[
\boxed{T(n)=O(n\log n)}.
\]

若选择使递推在 \(n=1\) 也成立的规范化初值，便可直接写成课件中的
\(T(1)/2+2(1/3+\cdots+1/(n+1))\)。若程序在长度为 1 时立即返回，从
\(n_0=2\) 开始展开即可，渐近结论一致。

调和级数产生对数，因为：

\[
\sum_{j=1}^{n}\frac1j\le 1+\int_1^n\frac{dx}{x}=1+\ln n.
\]

均衡递归树因为每次规模约减半而产生对数层数。这里则先对所有枢轴秩取平均，再由调和级数得到对数因子，因此允许部分划分很不均衡。

\hypertarget{ux4e3aux4ec0ux4e48ux968fux673aux7248ux672cux4ecdux53efux80fdux5f88ux6162}{%
\subsection{4.
为什么随机版本仍可能很慢}\label{ux4e3aux4ec0ux4e48ux968fux673aux7248ux672cux4ecdux53efux80fdux5f88ux6162}}

互异元素下，每次恰好选到当前最小值这一种具体序列的概率为：

\[
\frac1n\cdot\frac1{n-1}\cdots\frac12=\frac1{n!}.
\]

这个概率对应一条指定的极端路径。分析算法整体变慢的概率，还要考虑其他坏路径。上面的递推证明了期望为
\(O(n\log n)\)；若要得到指数衰减的尾界等更强的高概率结论，还需额外论证，不能由期望上界直接推出。

\hypertarget{ux8865ux5145ux63a8ux5bfcux7cbeux786eux6bd4ux8f83ux6b21ux6570}{%
\subsection{5.
补充推导：精确比较次数}\label{ux8865ux5145ux63a8ux5bfcux7cbeux786eux6bd4ux8f83ux6b21ux6570}}

如果只统计划分时的关键字比较，令 \(C_n\)
为期望比较次数，\(C_0=C_1=0\)，则：

\[
C_n=n-1+\frac2n\sum_{i=0}^{n-1}C_i.
\]

设排好序后各元素为 \(z_1<\cdots<z_n\)。一对 \(z_i,z_j\) 只有在区间
\(z_i,\ldots,z_j\) 中最先成为枢轴的是某个端点时才会相互比较，因此：

\[
\Pr(z_i\text{ 与 }z_j\text{ 比较})=\frac2{j-i+1}.
\]

利用期望的线性性求和：

\[
\mathbb E[C_n]
=\sum_{1\le i<j\le n}\frac2{j-i+1}
=2(n+1)H_n-4n,
\qquad H_n=\sum_{j=1}^{n}\frac1j.
\]

最坏比较次数为 \(n(n-1)/2\)。下图由这两个公式绘制，比较的是理论计数：

\begin{figure}
\centering
\includegraphics{../assets/2026-10-15-Day4/quicksort-comparisons.pdf}
\caption{互异元素下，随机快速排序的期望比较次数与最坏比较次数；双对数坐标展示增长速度。}
\end{figure}

期望比较次数约为 \(2n\ln n\)，最坏比较次数按二次增长。这里的 \(C_n\)
只统计关键字比较；前面将每层全部线性工作简写为 \(n\) 的
\(T(n)\)，采用的是另一种计量单位。

\hypertarget{ux516dquickselectux53eaux9012ux5f52ux5305ux542bux76eeux6807ux79e9ux7684ux4e00ux4fa7}{%
\section{六、Quickselect：只递归包含目标秩的一侧}\label{ux516dquickselectux53eaux9012ux5f52ux5305ux542bux76eeux6807ux79e9ux7684ux4e00ux4fa7}}

若目标是第 \(k\) 小元素，每次划分后只需判断枢轴的最终位置：

\begin{itemize}
\tightlist
\item
  枢轴位置恰好为目标位置，直接返回。
\item
  目标在左侧，只继续处理左侧。
\item
  目标在右侧，只继续处理右侧。
\end{itemize}

\Needspace{156pt}

\begin{Shaded}
\begin{Highlighting}[]
\KeywordTok{def}\NormalTok{ quickselect(a, rank, rng}\OperatorTok{=}\VariableTok{None}\NormalTok{):}
    \ControlFlowTok{if} \KeywordTok{not} \DecValTok{0} \OperatorTok{\textless{}=}\NormalTok{ rank }\OperatorTok{\textless{}} \BuiltInTok{len}\NormalTok{(a):}
        \ControlFlowTok{raise} \PreprocessorTok{ValueError}\NormalTok{(}\StringTok{"rank must be a valid zero{-}based index"}\NormalTok{)}
    \ControlFlowTok{if}\NormalTok{ rng }\KeywordTok{is} \VariableTok{None}\NormalTok{:}
\NormalTok{        rng }\OperatorTok{=}\NormalTok{ Random()}
\NormalTok{    p, r }\OperatorTok{=} \DecValTok{0}\NormalTok{, }\BuiltInTok{len}\NormalTok{(a) }\OperatorTok{{-}} \DecValTok{1}
    \ControlFlowTok{while}\NormalTok{ p }\OperatorTok{\textless{}=}\NormalTok{ r:}
\NormalTok{        q }\OperatorTok{=}\NormalTok{ randomized\_partition(a, p, r, rng)}
        \ControlFlowTok{if}\NormalTok{ q }\OperatorTok{==}\NormalTok{ rank:}
            \ControlFlowTok{return}\NormalTok{ a[q]}
        \ControlFlowTok{if}\NormalTok{ rank }\OperatorTok{\textless{}}\NormalTok{ q:}
\NormalTok{            r }\OperatorTok{=}\NormalTok{ q }\OperatorTok{{-}} \DecValTok{1}
        \ControlFlowTok{else}\NormalTok{:}
\NormalTok{            p }\OperatorTok{=}\NormalTok{ q }\OperatorTok{+} \DecValTok{1}
\end{Highlighting}
\end{Shaded}

这里的 \texttt{rank} 是 0-based 排名，例如奇数长度数组的中位数对应
\texttt{rank\ =\ len(a)\ //\ 2}。算法会改动输入数组。

若每次保留的规模约减半，总扫描代价为
\(n+n/2+n/4+\cdots=O(n)\)。随机枢轴下的期望线性时间，可以用归纳法证明。设
\(S(n)\) 为所有目标秩中最大的期望代价。固定任意目标秩，枢轴左边有 \(i\)
个元素时，记下一步实际处理的规模为 \(m_i\)；若恰好找到目标，则记为
0。于是：

\[
0\le m_i\le\max(i,n-1-i)<n.
\]

对于 \(n\ge2\)，

\[
\frac1n\sum_{i=0}^{n-1}\max(i,n-1-i)\le\frac{3n}{4}.
\]

用归纳假设 \(S(m)\le am\) 控制所有更小规模，对任何目标秩都有：

\[
\mathbb E[T]\le cn+\frac an\sum_{i=0}^{n-1}m_i
\le cn+\frac{3an}{4}\le an,
\qquad a\ge4c.
\]

同一个常数 \(a\) 对所有目标秩都成立，故
\(S(n)\le an\)，有限初值并入常数。因此，互异元素、均匀随机枢轴下的期望时间为
\(\Theta(n)\)，最坏仍为 \(\Theta(n^2)\)。

课件第 19 页将找中位数的期望时间写成 \(O(\log n)\)，应改为
\(O(n)\)。初始划分就要扫描线性数量的元素，分析总时间还需累加后续各层的扫描工作。要保证确定性的最坏线性时间，则需采用其他枢轴策略，例如
median-of-medians。

\hypertarget{ux4e03ux5806ux6392ux5e8fux4e0eux81eaux5e95ux5411ux4e0aux5efaux5806}{%
\section{七、堆排序与自底向上建堆}\label{ux4e03ux5806ux6392ux5e8fux4e0eux81eaux5e95ux5411ux4e0aux5efaux5806}}

\hypertarget{ux5efaux5806ux67e5ux770bux5806ux9876ux4e0eux53d6ux51faux5806ux9876}{%
\subsection{1.
建堆、查看堆顶与取出堆顶}\label{ux5efaux5806ux67e5ux770bux5806ux9876ux4e0eux53d6ux51faux5806ux9876}}

每次扫描剩余数组找最小值，连续取出 \(n\) 次的总代价为
\(\Theta(n^2)\)。最小堆要求父节点不大于子节点，因此全局最小值始终在堆顶。

\begin{longtable}[]{@{}
  >{\raggedright\arraybackslash}p{(\columnwidth - 4\tabcolsep) * \real{0.28}}
  >{\raggedright\arraybackslash}p{(\columnwidth - 4\tabcolsep) * \real{0.20}}
  >{\raggedright\arraybackslash}p{(\columnwidth - 4\tabcolsep) * \real{0.52}}@{}}
\toprule
操作 & 代价 & 做了什么 \\
\midrule
\endhead
自底向上建堆 & \(O(n)\) & 从整个输入建立堆序 \\
查看最小值 & \(O(1)\) & 读取堆顶，不删除 \\
取出最小值 & \(O(\log n)\) & 删除堆顶，并恢复堆序 \\
\bottomrule
\end{longtable}

课堂将最小堆视作黑盒：先建堆，再反复弹出最小值。

\Needspace{116pt}

\begin{Shaded}
\begin{Highlighting}[]
\ImportTok{import}\NormalTok{ heapq}


\KeywordTok{def}\NormalTok{ heap\_sort(a):}
\NormalTok{    heap }\OperatorTok{=} \BuiltInTok{list}\NormalTok{(a)}
\NormalTok{    heapq.heapify(heap)}
\NormalTok{    result }\OperatorTok{=}\NormalTok{ []}
    \ControlFlowTok{while}\NormalTok{ heap:}
\NormalTok{        result.append(heapq.heappop(heap))}
    \ControlFlowTok{return}\NormalTok{ result}
\end{Highlighting}
\end{Shaded}

建堆花费 \(O(n)\)，随后 \(n\) 次删除各需至多 \(O(\log n)\)，总时间为
\(O(n\log n)\)。上面的实现复制了输入，并单独保存输出，额外空间为
\(O(n)\)。教材中常见的 \textbf{原地堆排序}
使用最大堆，将最大值交换到数组末尾，辅助空间可降至
\(O(1)\)。空间结论取决于采用哪种实现。

\hypertarget{ux8865ux5145ux63a8ux5bfcux5efaux5806ux7684-onlog-n-ux4e0aux754cux4e3aux4f55ux53efux4ee5ux6536ux7d27}{%
\subsection{\texorpdfstring{2. 补充推导：建堆的 \(O(n\log n)\)
上界为何可以收紧}{2. 补充推导：建堆的 O(n\textbackslash log n) 上界为何可以收紧}}\label{ux8865ux5145ux63a8ux5bfcux5efaux5806ux7684-onlog-n-ux4e0aux754cux4e3aux4f55ux53efux4ee5ux6536ux7d27}}

逐个插入 \(n\) 个数，时间上界为 \(O(n\log n)\)。\texttt{heapify}
采用自底向上调整：大多数节点靠近叶子，调整距离很短，只有少数上层节点可能下降很多层。

高度为 \(h\) 的节点数量至多为 \(O(n/2^{h+1})\)，每个节点的调整代价为
\(O(h)\)。因此：

\[
\sum_{h\ge0}O\left(\frac{n}{2^{h+1}}\right)O(h)
=O(n)\sum_{h\ge0}\frac{h}{2^{h+1}}
=O(n).
\]

课堂主要使用建堆、取堆顶的接口及其复杂度。上面的求和补充了线性建堆的原因，后面的多路归并也会用到这一结论。

\hypertarget{ux516bux8ba1ux6570ux6392ux5e8fux4e0eux503cux57dfux5927ux5c0f}{%
\section{八、计数排序与值域大小}\label{ux516bux8ba1ux6570ux6392ux5e8fux4e0eux503cux57dfux5927ux5c0f}}

\hypertarget{ux7b97ux6cd5ux4e0eux4f8bux5b50}{%
\subsection{1. 算法与例子}\label{ux7b97ux6cd5ux4e0eux4f8bux5b50}}

假设输入都是正整数，值域为 \(1..K\)。建立计数数组 \(C[0..K]\)，让
\(C[x]\) 记录 \(x\) 出现的次数，再按 \(1,2,\ldots,K\) 的顺序输出。

对课堂例子 \([1,5,3,2,2]\)：

\begin{longtable}[]{@{}
  >{\raggedright\arraybackslash}p{(\columnwidth - 10\tabcolsep) * \real{0.30}}
  >{\raggedleft\arraybackslash}p{(\columnwidth - 10\tabcolsep) * \real{0.14}}
  >{\raggedleft\arraybackslash}p{(\columnwidth - 10\tabcolsep) * \real{0.14}}
  >{\raggedleft\arraybackslash}p{(\columnwidth - 10\tabcolsep) * \real{0.14}}
  >{\raggedleft\arraybackslash}p{(\columnwidth - 10\tabcolsep) * \real{0.14}}
  >{\raggedleft\arraybackslash}p{(\columnwidth - 10\tabcolsep) * \real{0.14}}@{}}
\toprule
值 \(x\) & 1 & 2 & 3 & 4 & 5 \\
\midrule
\endhead
次数 \(C[x]\) & 1 & 2 & 1 & 0 & 1 \\
\bottomrule
\end{longtable}

输出为 \([1,2,2,3,5]\)。

\Needspace{126pt}

\begin{Shaded}
\begin{Highlighting}[]
\KeywordTok{def}\NormalTok{ counting\_sort(a):}
    \ControlFlowTok{if} \KeywordTok{not}\NormalTok{ a:}
        \ControlFlowTok{return}\NormalTok{ []}
\NormalTok{    upper }\OperatorTok{=} \BuiltInTok{max}\NormalTok{(a)}
\NormalTok{    counts }\OperatorTok{=}\NormalTok{ [}\DecValTok{0}\NormalTok{] }\OperatorTok{*}\NormalTok{ (upper }\OperatorTok{+} \DecValTok{1}\NormalTok{)}
    \ControlFlowTok{for}\NormalTok{ x }\KeywordTok{in}\NormalTok{ a:}
\NormalTok{        counts[x] }\OperatorTok{+=} \DecValTok{1}
\NormalTok{    result }\OperatorTok{=}\NormalTok{ []}
    \ControlFlowTok{for}\NormalTok{ x }\KeywordTok{in} \BuiltInTok{range}\NormalTok{(}\DecValTok{1}\NormalTok{, upper }\OperatorTok{+} \DecValTok{1}\NormalTok{):}
\NormalTok{        result.extend([x] }\OperatorTok{*}\NormalTok{ counts[x])}
    \ControlFlowTok{return}\NormalTok{ result}
\end{Highlighting}
\end{Shaded}

这段代码要求输入为正整数；若要支持 0
或负整数，需要调整扫描范围或增加偏移。找最大值、统计频次及输出共需
\(O(n)\)，初始化和遍历值域需 \(O(K)\)，所以：

\[
\text{时间}=O(n+K),\qquad
\text{辅助空间}=O(K).
\]

若把返回的输出数组也计入额外存储，总空间为 \(O(n+K)\)。

这个简化版本只输出整数键值。若每条记录还带有其他字段，稳定排序就要保留完整记录及相同键的原有先后顺序，应采用累计计数与稳定放置等标准写法。重复输出相同整数，并不能完成这种记录排序。

\hypertarget{ux6570ux503cux5927ux5c0fux4e0eux4e8cux8fdbux5236ux7f16ux7801ux957fux5ea6}{%
\subsection{2.
数值大小与二进制编码长度}\label{ux6570ux503cux5927ux5c0fux4e0eux4e8cux8fdbux5236ux7f16ux7801ux957fux5ea6}}

复杂度式中的 \(K\)
表示数值大小，而计算复杂性中的「多项式时间」以输入编码长度 \(L\)
为参照。

正整数 \(a_i\) 的二进制长度是
\(\lfloor\lg a_i\rfloor+1\)。变长编码下，总输入长度还要加上分隔或长度信息；通常可以把主要部分写为：

\[
L\asymp\sum_{i=1}^{n}
\bigl(\lfloor\lg a_i\rfloor+1\bigr).
\]

若统一用足够容纳 \(1..K\) 的定长字段，单个字段需要
\(\lceil\lg(K+1)\rceil\) 位，总长度约为 \(n\lg K\)。

数值可以相对编码长度呈指数增长。例如数组只含一个数 \(2^b\)，存储它需要
\(b+1\) 位；若算法初始化从 1 到 \(2^b\) 的计数表，单这一步就需要
\(\Theta(2^b)\) 的时间和空间。

因此，\(O(n+K)\)
虽然是数值参数的多项式，却未必是二进制输入长度的多项式。这类依赖数值大小的界称为
\textbf{伪多项式时间}。若输入额外满足 \(K=O(n^c)\)，其中 \(c\)
为固定常数，则该算法在这一受限输入族上是多项式时间；当 \(K=O(n)\)
时，时间进一步降为线性。

\hypertarget{ux6876ux6392ux5e8fux57faux6570ux6392ux5e8fux4e0e-sleep-sort-ux7684ux6761ux4ef6}{%
\subsection{3. 桶排序、基数排序与 Sleep Sort
的条件}\label{ux6876ux6392ux5e8fux57faux6570ux6392ux5e8fux4e0e-sleep-sort-ux7684ux6761ux4ef6}}

课堂还提到了桶排序和基数排序，作为课件延伸。桶排序的期望线性时间依赖适当的输入分布和桶内处理条件；输入不是正整数这一点，本身不足以保证桶排序线性完成。基数排序则需计入位数及每一位的排序成本。

Sleep Sort
按数值设置等待时间，再按唤醒顺序输出。分析它时，要把计时、调度、并发资源和值域大小纳入计算模型；只数代码中的显式比较会漏掉这些代价。系统调度也未必能让任意输入都可靠有序。这个例子说明，复杂度结论必须与计算模型一起理解。

\hypertarget{ux4e5dux6bd4ux8f83ux6392ux5e8fux4e0bux754cux4e3aux4ec0ux4e48ux81f3ux5c11ux9700ux8981-omeganlog-n-ux6b21ux6bd4ux8f83}{%
\section{\texorpdfstring{九、比较排序下界：为什么至少需要
\(\Omega(n\log n)\)
次比较}{九、比较排序下界：为什么至少需要 \textbackslash Omega(n\textbackslash log n) 次比较}}\label{ux4e5dux6bd4ux8f83ux6392ux5e8fux4e0bux754cux4e3aux4ec0ux4e48ux81f3ux5c11ux9700ux8981-omeganlog-n-ux6b21ux6bd4ux8f83}}

\hypertarget{ux7528ux51b3ux7b56ux6811ux8868ux793aux7b97ux6cd5}{%
\subsection{1.
用决策树表示算法}\label{ux7528ux51b3ux7b56ux6811ux8868ux793aux7b97ux6cd5}}

限制算法只能通过比较获得键之间的顺序信息。对 \(n\)
个互异元素，一次比较只会产生两个可能结果，可表示为二叉树的一个分支。

树根对应第一次比较；向左或向右表示比较结果；后续节点对应由此前结果决定的下一次比较。到达叶子时，算法已经能确定一个正确的输出下标顺序。

一条根到叶子的路径长度，就是该输入执行的比较次数；树高就是最坏输入所需的比较次数。

\hypertarget{ux4e3aux4ec0ux4e48ux81f3ux5c11ux8981ux6709-n-ux4e2aux53f6ux5b50}{%
\subsection{\texorpdfstring{2. 为什么至少要有 \(n!\)
个叶子}{2. 为什么至少要有 n! 个叶子}}\label{ux4e3aux4ec0ux4e48ux81f3ux5c11ux8981ux6709-n-ux4e2aux53f6ux5b50}}

把输入位置标记为 \(a_1,\ldots,a_n\)。当键互异时，它们可能有 \(n!\)
种严格全序。一个固定的输出下标顺序，例如
\((a_2,a_1,a_3)\)，只能在与它匹配的全序下正确。

正确的比较排序必须覆盖所有这些可能性。两个需要不同输出顺序的全序不能最终共享同一个输出叶子，因此叶子数至少是
\(n!\)。

这里区分的是元素相对顺序。即使键可以取无限多种数值，只要全序相同，它们仍可以沿同一条比较路径到达同一叶子。

\hypertarget{ux7531ux53f6ux5b50ux6570ux63a8ux51faux9ad8ux5ea6}{%
\subsection{3.
由叶子数推出高度}\label{ux7531ux53f6ux5b50ux6570ux63a8ux51faux9ad8ux5ea6}}

高度为 \(h\) 的二叉树最多有 \(2^h\) 个叶子，因此：

\[
2^h\ge n!,
\qquad
h\ge\lceil\lg(n!)\rceil.
\]

再证明 \(\lg(n!)=\Theta(n\lg n)\)。上界直接用 \(n!\le n^n\)：

\[
\lg(n!)\le n\lg n.
\]

下界取阶乘中最大的至少 \(n/2\) 个因子，每个都不小于 \(n/2\)：

\[
n!\ge (n/2)^{n/2},
\qquad
\lg(n!)\ge\frac n2\lg\frac n2=\Omega(n\lg n).
\]

由此得到：

\[
\boxed{\text{任意正确的比较排序，最坏情况需要 }
\Omega(n\log n)\text{ 次比较。}}
\]

课件第 30 页在阶乘附近标注了 \(\Theta(n^n)\)。这里应使用粗上界
\(n!\le n^n\)，并在取对数后得到
\(\lg(n!)=\Theta(n\lg n)\)，不能把前一个标记当作阶乘的等价式。

例如 \(n=3\) 时有 6 种全序，二叉树至少要有 6 个叶子，最坏路径至少为
\(\lceil\lg6\rceil=3\)。一般的叶子计数只给出下界，不保证存在某个算法恰好达到这个整数值。

\hypertarget{ux4e0bux754cux7684ux9002ux7528ux8303ux56f4}{%
\subsection{4.
下界的适用范围}\label{ux4e0bux754cux7684ux9002ux7528ux8303ux56f4}}

该下界针对最坏情况。带有序性预检查的算法遇到已排序输入时，仍可在线性扫描后返回；某些输入较容易处理，与最坏情况下界一致。

计数排序通过键值直接寻址计数表，利用了比较模型以外的信息，因此其
\(O(n+K)\)
时间与比较排序下界相容。特定分布、重复值结构及其他受限输入族，也可能允许更快的算法。

推广到随机算法的期望下界，还需对输入分布取平均。让 \(n!\)
种全序等概率出现，任何固定的正确比较决策树，平均叶深都至少为
\(\lg(n!)\)。再对随机选择取平均，平均比较次数仍为
\(\Omega(n\log n)\)，因此至少有一个输入的期望比较次数达到这一量级。这里要区分两个量词顺序：固定随机选择后存在坏输入，并不等于存在一个输入对所有随机选择都坏。

一次二元比较至多提供一比特的分支信息，而排序需要区分 \(n!\)
种顺序。决策树用叶子数与树高的关系，给出了这一信息量直觉的严格证明。

\hypertarget{ux5341ux5e38ux89c1ux6392ux5e8fux7b97ux6cd5ux7684ux590dux6742ux5ea6ux5bf9ux7167}{%
\section{十、常见排序算法的复杂度对照}\label{ux5341ux5e38ux89c1ux6392ux5e8fux7b97ux6cd5ux7684ux590dux6742ux5ea6ux5bf9ux7167}}

下表列出常见实现的复杂度，以及平均或期望结论所需的条件。

\begin{longtable}[]{@{}
  >{\raggedright\arraybackslash}p{(\columnwidth - 6\tabcolsep) * \real{0.20}}
  >{\raggedright\arraybackslash}p{(\columnwidth - 6\tabcolsep) * \real{0.21}}
  >{\raggedright\arraybackslash}p{(\columnwidth - 6\tabcolsep) * \real{0.25}}
  >{\raggedright\arraybackslash}p{(\columnwidth - 6\tabcolsep) * \real{0.34}}@{}}
\toprule
算法 & 最坏时间 & 平均或期望时间 & 辅助空间与条件 \\
\midrule
\endhead
插入排序 & \(\Theta(n^2)\) & 随机排列下 \(\Theta(n^2)\) &
\(O(1)\)；近乎有序时为 \(\Theta(n+I)\) \\
归并排序 & \(\Theta(n\log n)\) & \(\Theta(n\log n)\) & 常见数组实现为
\(O(n)\) \\
随机快速排序 & \(\Theta(n^2)\) & 互异元素下 \(\Theta(n\log n)\) &
原地划分；常见递归栈期望 \(O(\log n)\) \\
原地堆排序 & \(O(n\log n)\) & \(O(n\log n)\) 上界 &
\(O(1)\)；使用最大堆原地放置 \\
计数排序 & \(\Theta(n+K)\) & \(\Theta(n+K)\) & 计数表 \(O(K)\)；含输出为
\(O(n+K)\) \\
\bottomrule
\end{longtable}

其中 \(I\) 为当前数组的逆序对数：若 \(i<j\) 且
\(A[i]>A[j]\)，两者就构成一个逆序对。插入排序每次向右移动一个前驱元素，都会消除一个逆序对，第三道作业将据此计算移动次数。

\hypertarget{ux5341ux4e00ux4f5cux4e1a-1ux975eux9012ux5f52ux7684ux7ebfux6027ux65f6ux95f4ux6700ux5927ux5b50ux6570ux7ec4}{%
\section{十一、作业
1：非递归的线性时间最大子数组}\label{ux5341ux4e00ux4f5cux4e1a-1ux975eux9012ux5f52ux7684ux7ebfux6027ux65f6ux95f4ux6700ux5927ux5b50ux6570ux7ec4}}

\hypertarget{ux9898ux610fux4e0eux72b6ux6001ux5206ux6790}{%
\subsection{1.
题意与状态分析}\label{ux9898ux610fux4e0eux72b6ux6001ux5206ux6790}}

来源：课件第 32 页，CLRS 练习 \textbf{4.1-5}。

给定非空数组 \(A\)，元素可以为正、负或零。求一个
\textbf{非空连续子数组}，使元素和最大。要求从左向右扫描，采用非递归算法，在线性时间内完成。

题目提示：已知前缀 \(A[0..j]\) 的答案，扩展到 \(A[0..j+1]\)
后，最优解只有两种位置关系：

\begin{enumerate}
\def\labelenumi{\arabic{enumi}.}
\tightlist
\item
  完全位于旧前缀中。
\item
  以新元素 \(A[j+1]\) 结尾。
\end{enumerate}

因此需要同时记录整个前缀的最佳值，以及恰好在当前位置结束的最佳值。前者保存已有答案，后者决定新元素接在哪一段后面；只保留全局最大值会缺少后一项信息。

\hypertarget{ux4e24ux4e2aux72b6ux6001ux4e0eux9012ux63a8}{%
\subsection{2.
两个状态与递推}\label{ux4e24ux4e2aux72b6ux6001ux4e0eux9012ux63a8}}

定义：

\[
E_j=\max_{0\le i\le j}\sum_{t=i}^{j}A[t],
\]

即以位置 \(j\) 结尾的最大非空子数组和。另设：

\[
B_j=\max_{0\le i\le t\le j}\sum_{u=i}^{t}A[u],
\]

即整个前缀 \(A[0..j]\) 中的最大非空子数组和。

以 \(j\) 结尾的最优区间，要么只选 \(A[j]\)，要么把一个以 \(j-1\)
结尾的最优区间延长：

\[
E_j=\max\{A[j],E_{j-1}+A[j]\}.
\]

前缀最佳值则为：

\[
B_j=\max\{B_{j-1},E_j\}.
\]

初值为
\(E_0=B_0=A[0]\)。旧后缀贡献为负时，从当前元素重新开始更优。这一选择取决于以当前位置结束的状态，不能由此前的全局最佳值决定。

\hypertarget{ux8fd4ux56deux6700ux5927ux548cux4e0eux5de6ux53f3ux7aefux70b9ux7684ux5b9eux73b0}{%
\subsection{3.
返回最大和与左右端点的实现}\label{ux8fd4ux56deux6700ux5927ux548cux4e0eux5de6ux53f3ux7aefux70b9ux7684ux5b9eux73b0}}

\Needspace{186pt}

\begin{Shaded}
\begin{Highlighting}[]
\KeywordTok{def}\NormalTok{ max\_subarray(a):}
    \ControlFlowTok{if} \KeywordTok{not}\NormalTok{ a:}
        \ControlFlowTok{raise} \PreprocessorTok{ValueError}\NormalTok{(}\StringTok{"nonempty subarray requires nonempty input"}\NormalTok{)}
\NormalTok{    end\_sum }\OperatorTok{=}\NormalTok{ best\_sum }\OperatorTok{=}\NormalTok{ a[}\DecValTok{0}\NormalTok{]}
\NormalTok{    end\_left }\OperatorTok{=}\NormalTok{ best\_left }\OperatorTok{=}\NormalTok{ best\_right }\OperatorTok{=} \DecValTok{0}

    \ControlFlowTok{for}\NormalTok{ j }\KeywordTok{in} \BuiltInTok{range}\NormalTok{(}\DecValTok{1}\NormalTok{, }\BuiltInTok{len}\NormalTok{(a)):}
        \ControlFlowTok{if}\NormalTok{ end\_sum }\OperatorTok{+}\NormalTok{ a[j] }\OperatorTok{\textless{}}\NormalTok{ a[j]:}
\NormalTok{            end\_sum }\OperatorTok{=}\NormalTok{ a[j]}
\NormalTok{            end\_left }\OperatorTok{=}\NormalTok{ j}
        \ControlFlowTok{else}\NormalTok{:}
\NormalTok{            end\_sum }\OperatorTok{+=}\NormalTok{ a[j]}
        \ControlFlowTok{if}\NormalTok{ end\_sum }\OperatorTok{\textgreater{}}\NormalTok{ best\_sum:}
\NormalTok{            best\_sum }\OperatorTok{=}\NormalTok{ end\_sum}
\NormalTok{            best\_left, best\_right }\OperatorTok{=}\NormalTok{ end\_left, j}

    \ControlFlowTok{return}\NormalTok{ best\_sum, best\_left, best\_right}
\end{Highlighting}
\end{Shaded}

返回值为
\texttt{(最大和,\ 左端点,\ 右端点)}。相等时保留已有答案，因此有多个最优区间时，会返回最先记录的那个，满足题目求一个最优区间的要求。

\hypertarget{ux624bux7b97ux5b8cux6574ux8fc7ux7a0b}{%
\subsection{4.
手算完整过程}\label{ux624bux7b97ux5b8cux6574ux8fc7ux7a0b}}

取：

\[
A=[-2,1,-3,4,-1,2,1,-5,4].
\]

\begin{longtable}[]{@{}
  >{\raggedright\arraybackslash}p{(\columnwidth - 10\tabcolsep) * \real{0.07}}
  >{\raggedleft\arraybackslash}p{(\columnwidth - 10\tabcolsep) * \real{0.10}}
  >{\raggedleft\arraybackslash}p{(\columnwidth - 10\tabcolsep) * \real{0.10}}
  >{\raggedright\arraybackslash}p{(\columnwidth - 10\tabcolsep) * \real{0.23}}
  >{\raggedleft\arraybackslash}p{(\columnwidth - 10\tabcolsep) * \real{0.13}}
  >{\raggedright\arraybackslash}p{(\columnwidth - 10\tabcolsep) * \real{0.37}}@{}}
\toprule
\(j\) & \(A[j]\) & \(E_j\) & 当前结尾区间 & \(B_j\) &
已记录的最佳区间 \\
\midrule
\endhead
0 & -2 & -2 & \([0,0]\) & -2 & \([0,0]\) \\
1 & 1 & 1 & \([1,1]\) & 1 & \([1,1]\) \\
2 & -3 & -2 & \([1,2]\) & 1 & \([1,1]\) \\
3 & 4 & 4 & \([3,3]\) & 4 & \([3,3]\) \\
4 & -1 & 3 & \([3,4]\) & 4 & \([3,3]\) \\
5 & 2 & 5 & \([3,5]\) & 5 & \([3,5]\) \\
6 & 1 & 6 & \([3,6]\) & 6 & \([3,6]\) \\
7 & -5 & 1 & \([3,7]\) & 6 & \([3,6]\) \\
8 & 4 & 5 & \([3,8]\) & 6 & \([3,6]\) \\
\bottomrule
\end{longtable}

答案为 \([4,-1,2,1]\)，总和为 6，端点为 \([3,6]\)。处理到位置 7
时，结尾状态的和下降，但记录的全局最佳区间保持不变。

\hypertarget{ux6b63ux786eux6027ux8bc1ux660e}{%
\subsection{5. 正确性证明}\label{ux6b63ux786eux6027ux8bc1ux660e}}

循环处理完下标 \(j\) 后，保持两个不变式：

\begin{enumerate}
\def\labelenumi{\arabic{enumi}.}
\tightlist
\item
  \texttt{end\_sum} 等于所有以 \(j\)
  结尾的非空子数组中的最大和，\texttt{end\_left} 是对应起点。
\item
  \texttt{best\_sum} 等于前缀 \(A[0..j]\)
  中所有非空子数组的最大和，记录的端点对应一个达到该和的区间。
\end{enumerate}

初始 \(j=0\) 时只有一个非空区间 \([0,0]\)，两个状态均取 \(A[0]\)，成立。

假设处理完 \(j-1\) 时两个不变式成立。以 \(j\) 结尾的非空区间，要么从
\(j\) 开始，要么由以 \(j-1\)
结尾的区间延长。后一类都加上同一个新元素，选旧区间中和最大的那个即可，因此第一次更新得到
\(E_j\)。前缀的最优区间则要么不经过 \(j\)，要么以 \(j\)
结尾，第二次更新据此得到
\(B_j\)。两次选择都同步维护端点，所以不变式继续成立。

循环结束时 \(j=n-1\)，前缀就是整个数组，因此返回值正确。

\hypertarget{ux8fb9ux754cux4e0eux590dux6742ux5ea6}{%
\subsection{6.
边界与复杂度}\label{ux8fb9ux754cux4e0eux590dux6742ux5ea6}}

数组全为负数时，例如 \([-5,-2,-8]\)，答案是单个元素 \([-2]\)，返回
\((-2,1,1)\)。若把最佳值初始化为 0，便会错误地返回一个不允许的空子数组。

长度为 1
时，直接返回唯一元素。空数组不满足输入要求，代码会抛出异常。若题目改为允许空子数组，可以将答案下限设为
0，但也要相应修改状态定义与证明，和本题的非空约定区分开。

每个元素只更新固定数量的变量，因此时间为 \(\Theta(n)\)，辅助空间为
\(O(1)\)。上一课的分治方法花费
\(\Theta(n\log n)\)；这里用一个最佳结尾状态，保存下一次扩展所需的跨边界信息。

\hypertarget{ux5341ux4e8cux4f5cux4e1a-2ux7528ux6700ux5c0fux5806ux5408ux5e76-k-ux4e2aux6709ux5e8fux8868}{%
\section{\texorpdfstring{十二、作业 2：用最小堆合并 \(k\)
个有序表}{十二、作业 2：用最小堆合并 k 个有序表}}\label{ux5341ux4e8cux4f5cux4e1a-2ux7528ux6700ux5c0fux5806ux5408ux5e76-k-ux4e2aux6709ux5e8fux8868}}

\hypertarget{ux9898ux610fux4e0eux5173ux952eux89c2ux5bdf}{%
\subsection{1.
题意与关键观察}\label{ux9898ux610fux4e0eux5173ux952eux89c2ux5bdf}}

来源：课件第 33 页，CLRS 练习 \textbf{6.5-9}。

给定 \(k\) 个已经分别按非降序排列的表，总元素数为
\(n\)。将它们合成一个有序表，要求时间
\(O(n\lg k)\)。提示是使用最小堆进行多路归并。

每张有序表中，第一个尚未输出的元素就是剩余部分的最小值。因此，只需比较各非空表的当前表头，就能找到下一个输出值。

堆中保存至多 \(k\) 个候选值。取出一个表头后，只需从它所属的表补入后继。

\hypertarget{ux5806ux8282ux70b9ux4fddux5b58ux4ec0ux4e48}{%
\subsection{2.
堆节点保存什么}\label{ux5806ux8282ux70b9ux4fddux5b58ux4ec0ux4e48}}

对于数组形式的输入表，堆节点设为：

\[
(\text{value},\text{list\_id},\text{index}).
\]

第一个字段决定数值大小；后两个字段记录来源，便于找到后继。值相等时，Python
元组继续比较表号和下标，因此可以正常处理重复值，也不需要比较不可排序的自定义节点。

\Needspace{176pt}

\begin{Shaded}
\begin{Highlighting}[]
\KeywordTok{def}\NormalTok{ merge\_sorted\_lists(lists):}
\NormalTok{    heap }\OperatorTok{=}\NormalTok{ [}
\NormalTok{        (row[}\DecValTok{0}\NormalTok{], list\_id, }\DecValTok{0}\NormalTok{)}
        \ControlFlowTok{for}\NormalTok{ list\_id, row }\KeywordTok{in} \BuiltInTok{enumerate}\NormalTok{(lists)}
        \ControlFlowTok{if}\NormalTok{ row}
\NormalTok{    ]}
\NormalTok{    heapq.heapify(heap)}
\NormalTok{    result }\OperatorTok{=}\NormalTok{ []}
    \ControlFlowTok{while}\NormalTok{ heap:}
\NormalTok{        value, list\_id, index }\OperatorTok{=}\NormalTok{ heapq.heappop(heap)}
\NormalTok{        result.append(value)}
\NormalTok{        next\_index }\OperatorTok{=}\NormalTok{ index }\OperatorTok{+} \DecValTok{1}
        \ControlFlowTok{if}\NormalTok{ next\_index }\OperatorTok{\textless{}} \BuiltInTok{len}\NormalTok{(lists[list\_id]):}
\NormalTok{            next\_value }\OperatorTok{=}\NormalTok{ lists[list\_id][next\_index]}
\NormalTok{            heapq.heappush(heap, (next\_value, list\_id, next\_index))}
    \ControlFlowTok{return}\NormalTok{ result}
\end{Highlighting}
\end{Shaded}

若输入是链表，则每个堆节点保存当前节点及表号，弹出后沿 \texttt{next}
指针加入后继，算法结构不变。

\hypertarget{ux4e00ux6b21ux5408ux5e76ux7684ux5168ux8fc7ux7a0b}{%
\subsection{3.
一次合并的全过程}\label{ux4e00ux6b21ux5408ux5e76ux7684ux5168ux8fc7ux7a0b}}

设：

\[
L_0=[1,4,7],\quad L_1=[1,3,8],\quad L_2=[2,6].
\]

初始候选为
\((1,0,0),(1,1,0),(2,2,0)\)。下表为便于阅读，按数值和表号列出候选集合；实际堆数组的排列可能不同：

\begin{longtable}[]{@{}
  >{\raggedright\arraybackslash}p{(\columnwidth - 6\tabcolsep) * \real{0.13}}
  >{\raggedright\arraybackslash}p{(\columnwidth - 6\tabcolsep) * \real{0.25}}
  >{\raggedright\arraybackslash}p{(\columnwidth - 6\tabcolsep) * \real{0.18}}
  >{\raggedright\arraybackslash}p{(\columnwidth - 6\tabcolsep) * \real{0.44}}@{}}
\toprule
第几次取出 & 输出值与来源 & 加入的后继 & 取出后尚在堆中的候选值 \\
\midrule
\endhead
1 & 1，来自 \(L_0\) & 4 & \(1(L_1),2(L_2),4(L_0)\) \\
2 & 1，来自 \(L_1\) & 3 & \(2(L_2),3(L_1),4(L_0)\) \\
3 & 2，来自 \(L_2\) & 6 & \(3(L_1),4(L_0),6(L_2)\) \\
4 & 3，来自 \(L_1\) & 8 & \(4(L_0),6(L_2),8(L_1)\) \\
5 & 4，来自 \(L_0\) & 7 & \(6(L_2),7(L_0),8(L_1)\) \\
6 & 6，来自 \(L_2\) & 无 & \(7(L_0),8(L_1)\) \\
7 & 7，来自 \(L_0\) & 无 & \(8(L_1)\) \\
8 & 8，来自 \(L_1\) & 无 & 空 \\
\bottomrule
\end{longtable}

最终结果是
\([1,1,2,3,4,6,7,8]\)。每个输入元素进入输出一次；表耗尽后不再向堆中放入候选。

\hypertarget{ux6b63ux786eux6027ux8bc1ux660e-1}{%
\subsection{4. 正确性证明}\label{ux6b63ux786eux6027ux8bc1ux660e-1}}

循环开始时保持：

\begin{enumerate}
\def\labelenumi{\arabic{enumi}.}
\tightlist
\item
  对每个仍有元素未输出的表，堆中恰有它的第一个未输出元素。
\item
  已输出部分有序，并且不大于任何尚未输出的元素。
\end{enumerate}

初始化把各非空表的第一个元素加入堆，输出为空，成立。

各表内部有序，所以未输出元素都不小于本表的当前表头。堆顶又是所有当前表头中的最小值，也就是全体未输出元素的最小值。将它接在输出末尾，仍保持有序；再加入同一张表的后继，便恢复了每张未耗尽的表恰有一个候选值的不变式。

每轮恰好输出一个元素，循环共进行 \(n\)
轮。堆空时，各表都已经耗尽，所以输出有序并且包含全部输入元素。

\hypertarget{ux65f6ux95f4ux7a7aux95f4ux4e0eux7279ux6b8aux60c5ux51b5}{%
\subsection{5.
时间、空间与特殊情况}\label{ux65f6ux95f4ux7a7aux95f4ux4e0eux7279ux6b8aux60c5ux51b5}}

若 \(k\) 个表都非空且 \(k\ge2\)，初始化用 \(O(k)\)
时间，\texttt{heapify} 也是 \(O(k)\)；共弹出 \(n\)
次，每次至多再插入一次，堆大小始终不超过 \(k\)，因此：

\[
T(n,k)=O(k+n\lg k)=O(n\lg k),
\]

最后一步使用了 \(k\le n\)。堆的辅助空间为
\(O(k)\)，若把返回的输出也计入存储，则为 \(O(n+k)\)。

如果允许大量空表，扫描这 \(k\) 个输入表头本身就需
\(O(k)\)。设其中非空表数为 \(k'\)，更完整的表达是：

\[
O\bigl(k+n\log\max(2,k')\bigr).
\]

\(k=0\) 时返回空表；只有一个非空表时，复制输出仍需 \(\Theta(n)\)，即使
\(\log1=0\)，总时间也不为零。题目常用的 \(O(n\lg k)\)
写法默认至少合并两张非空表。

逐个扫描所有当前表头，每输出一个元素需要 \(O(k)\)，总计
\(O(nk)\)；最小堆把这一步的维护成本降到了 \(O(\log k)\)。

\hypertarget{ux5341ux4e09ux4f5cux4e1a-3ux622aux65adux5febux901fux6392ux5e8fux518dux5bf9ux5168ux6570ux7ec4ux505aux63d2ux5165ux6392ux5e8f}{%
\section{十三、作业
3：截断快速排序，再对全数组做插入排序}\label{ux5341ux4e09ux4f5cux4e1a-3ux622aux65adux5febux901fux6392ux5e8fux518dux5bf9ux5168ux6570ux7ec4ux505aux63d2ux5165ux6392ux5e8f}}

\hypertarget{ux9898ux610fux4e2dux7684ux6267ux884cux987aux5e8f}{%
\subsection{1.
题意中的执行顺序}\label{ux9898ux610fux4e2dux7684ux6267ux884cux987aux5e8f}}

来源：课件第 34 页，CLRS 练习 \textbf{7.4-5}。

修改快速排序：调用到一个 \textbf{少于 \(k\) 个元素}
的子数组时，直接返回，让这个小区间暂时无序。整个快速排序阶段结束后，再对
\textbf{整个数组} 执行一次插入排序。

要求证明期望时间为：

\[
\boxed{O\bigl(nk+n\lg(n/k)\bigr)},
\]

并讨论理论上、实践中如何选择 \(k\)。

分析采用互异元素和均匀随机枢轴，设 \(2\le k\le n\)。截止条件严格为
\(m<k\)，长度恰好为 \(k\) 的子数组仍要划分。

\hypertarget{ux5b8cux6574ux7b97ux6cd5}{%
\subsection{2. 完整算法}\label{ux5b8cux6574ux7b97ux6cd5}}

\Needspace{246pt}

\begin{Shaded}
\begin{Highlighting}[]
\KeywordTok{def}\NormalTok{ hybrid\_sort(a, k, rng}\OperatorTok{=}\VariableTok{None}\NormalTok{):}
    \ControlFlowTok{if}\NormalTok{ k }\OperatorTok{\textless{}} \DecValTok{2}\NormalTok{:}
        \ControlFlowTok{raise} \PreprocessorTok{ValueError}\NormalTok{(}\StringTok{"cutoff k must be at least 2"}\NormalTok{)}
    \ControlFlowTok{if}\NormalTok{ rng }\KeywordTok{is} \VariableTok{None}\NormalTok{:}
\NormalTok{        rng }\OperatorTok{=}\NormalTok{ Random()}

    \KeywordTok{def}\NormalTok{ partial\_quicksort(p, r):}
        \ControlFlowTok{if}\NormalTok{ r }\OperatorTok{{-}}\NormalTok{ p }\OperatorTok{+} \DecValTok{1} \OperatorTok{\textless{}}\NormalTok{ k:}
            \ControlFlowTok{return}
\NormalTok{        q }\OperatorTok{=}\NormalTok{ randomized\_partition(a, p, r, rng)}
\NormalTok{        partial\_quicksort(p, q }\OperatorTok{{-}} \DecValTok{1}\NormalTok{)}
\NormalTok{        partial\_quicksort(q }\OperatorTok{+} \DecValTok{1}\NormalTok{, r)}

\NormalTok{    partial\_quicksort(}\DecValTok{0}\NormalTok{, }\BuiltInTok{len}\NormalTok{(a) }\OperatorTok{{-}} \DecValTok{1}\NormalTok{)}

    \ControlFlowTok{for}\NormalTok{ j }\KeywordTok{in} \BuiltInTok{range}\NormalTok{(}\DecValTok{1}\NormalTok{, }\BuiltInTok{len}\NormalTok{(a)):}
\NormalTok{        key }\OperatorTok{=}\NormalTok{ a[j]}
\NormalTok{        i }\OperatorTok{=}\NormalTok{ j }\OperatorTok{{-}} \DecValTok{1}
        \ControlFlowTok{while}\NormalTok{ i }\OperatorTok{\textgreater{}=} \DecValTok{0} \KeywordTok{and}\NormalTok{ a[i] }\OperatorTok{\textgreater{}}\NormalTok{ key:}
\NormalTok{            a[i }\OperatorTok{+} \DecValTok{1}\NormalTok{] }\OperatorTok{=}\NormalTok{ a[i]}
\NormalTok{            i }\OperatorTok{{-}=} \DecValTok{1}
\NormalTok{        a[i }\OperatorTok{+} \DecValTok{1}\NormalTok{] }\OperatorTok{=}\NormalTok{ key}
    \ControlFlowTok{return}\NormalTok{ a}
\end{Highlighting}
\end{Shaded}

代码复用前面的 \texttt{randomized\_partition}，可接受空数组及 \(k>n\)
的情况；后者会直接进入全数组插入排序。题目中含 \(\lg(n/k)\)
的复杂度分析，采用的范围为 \(2\le k\le n\)。

前面的划分只交换元素，最后的插入排序又会排好整个数组，因此算法正确。分析运行时间时，需要先说明第一阶段留下的数组有多接近有序。

\hypertarget{ux672aux6392ux5e8fux5c0fux5757ux4e4bux95f4ux7684ux987aux5e8f}{%
\subsection{3.
未排序小块之间的顺序}\label{ux672aux6392ux5e8fux5c0fux5757ux4e4bux95f4ux7684ux987aux5e8f}}

递归在小于 \(k\)
的区间上停止。这些终止区间内部可以无序，但它们之间的顺序已经由祖先枢轴确定。

任取两个不同的终止区间，考虑把它们首次分开的祖先划分。左侧全部元素小于该枢轴，右侧全部元素大于该枢轴，所以左区间中的每个元素都小于右区间中的每个元素。已经就位的枢轴也不会与其他区间形成逆序。

因此，剩余逆序对只可能出现在同一个未排序小块内部。

以 \(k=4\) 为例，一种合法的第一阶段结果是：

\[
[3,1,2]\;|\;4\;|\;[6,5]\;|\;7\;|\;[9,8].
\]

其中 4 和 7 是已经固定的枢轴，三个小块的长度分别为 3、2、2，均小于
4。小块内部尚未排序，块间已满足前块所有值小于后块所有值的关系。

\hypertarget{ux63d2ux5165ux6392ux5e8fux7684ux4ee3ux4ef7ux7528ux9006ux5e8fux5bf9ux8ba1ux91cf}{%
\subsection{4.
插入排序的代价用逆序对计量}\label{ux63d2ux5165ux6392ux5e8fux7684ux4ee3ux4ef7ux7528ux9006ux5e8fux5bf9ux8ba1ux91cf}}

标准插入排序插入 \texttt{key}
时，每向右移动一个比它大的前驱元素，就消除一个逆序对。每个原始逆序对恰好引起一次这种移动，加上每轮固定开销，运行时间为：

\[
\Theta(n+I),
\]

其中 \(I\) 是第一阶段结束后的逆序对数。

设未排序小块长度为 \(m_1,m_2,\ldots,m_s\)。这些块互不重叠，满足
\(m_i<k\)、\(\sum_i m_i\le n\)。长度为 \(m_i\) 的块至多有
\(\binom{m_i}{2}\) 个逆序对，所以：

\[
I\le\sum_{i=1}^{s}\binom{m_i}{2}
\le\frac{k-1}{2}\sum_{i=1}^{s}m_i
\le\frac{n(k-1)}2
=O(nk).
\]

因此，最后一次 \textbf{全数组} 插入排序需要
\(O(n+nk)=O(nk)\)。全局逆序对数限制了整个扫描过程的移动次数，所以即使不逐块排序，也能得到这个上界。

前面的具体例子中，\([3,1,2]\) 有 2 个逆序对，\([6,5]\) 和 \([9,8]\) 各有
1 个，总计 \(I=4\)。完整数组的插入过程只需 4 次右移：

\begin{longtable}[]{@{}
  >{\raggedright\arraybackslash}p{(\columnwidth - 4\tabcolsep) * \real{0.20}}
  >{\raggedright\arraybackslash}p{(\columnwidth - 4\tabcolsep) * \real{0.30}}
  >{\raggedright\arraybackslash}p{(\columnwidth - 4\tabcolsep) * \real{0.50}}@{}}
\toprule
插入的键 & 被右移的元素 & 局部结果 \\
\midrule
\endhead
1 & 3 & \([1,3,2]\) \\
2 & 3 & \([1,2,3]\) \\
5 & 6 & \([5,6]\) \\
8 & 9 & \([8,9]\) \\
\bottomrule
\end{longtable}

其余插入轮次只进行固定数量的检查。最终数组为 \([1,2,3,4,5,6,7,8,9]\)。

\hypertarget{ux622aux65adux5febux901fux6392ux5e8fux7684ux671fux671bux4ee3ux4ef7}{%
\subsection{5.
截断快速排序的期望代价}\label{ux622aux65adux5febux901fux6392ux5e8fux7684ux671fux671bux4ee3ux4ef7}}

若每次恰好均分，递归约有 \(\lg(n/k)\) 层，每层处理 \(O(n)\)
个元素，于是得到
\(O(n\lg(n/k))\)。这是均衡划分下的直观分析。要覆盖随机快速排序中的不均衡划分，还需计算达到截止规模之前的期望工作量。

将长度为 \(m\) 的一次划分所需的 \(\Theta(m)\)
工作，平均分摊给区间内的每个元素。固定任意原始元素
\(x\)，跟踪它在尚需划分的子问题中经历的区间规模。

称某次枢轴为「好枢轴」，若它的 1-based 秩 \(r\) 满足：

\[
\lfloor m/4\rfloor+1\le r\le\lceil3m/4\rceil.
\]

这些秩至少占一半，所以每次条件概率至少为 \(1/2\)。取到好枢轴后，若 \(x\)
本身就是枢轴，它退出递归；否则它进入的那一侧规模至多为
\(3m/4\)，取整不会破坏这个上界。

从规模 \(n\) 缩小到小于 \(k\)，至多需要：

\[
g=\left\lfloor\log_{4/3}\frac nk\right\rfloor+1
\]

次这样的成功缩小。即使中间出现坏划分，每等到一次好枢轴，期望尝试次数也至多为
2：给定此前任意划分历史，下一次选中好枢轴的条件概率始终至少为
\(1/2\)。如果元素提前成为枢轴，跟踪还会更早结束。

因此，一个元素参与划分的期望次数为：

\[
O\left(1+\log\frac nk\right).
\]

再对 \(n\) 个元素求和，利用期望的线性性：

\[
\mathbb E[\text{第一阶段代价}]
=O\left(n+n\log\frac nk\right).
\]

不同元素的递归路径可以相互依赖。上述求和只用到每步随机枢轴的条件概率，以及期望的线性性，无须假设路径独立或每层均衡。

\hypertarget{ux5408ux5e76ux4e24ux9636ux6bb5ux7684ux754c}{%
\subsection{6.
合并两阶段的界}\label{ux5408ux5e76ux4e24ux9636ux6bb5ux7684ux754c}}

加入插入排序以及递归调用的线性级开销：

\[
\mathbb E[T(n,k)]
=O\left(n+n\lg\frac nk+nk\right)
=\boxed{O\left(nk+n\lg\frac nk\right)}.
\]

最后一步利用 \(k\ge2\)，将 \(O(n)\) 吸收到 \(O(nk)\) 中。特别地，\(k=n\)
时虽然 \(\lg(n/k)=0\)，第一阶段仍需一次线性划分。这部分由总界中的 \(nk\)
项覆盖，单独计算第一阶段时仍要保留。

这个界是随机枢轴下的期望结论。若连续选到很差的枢轴，第一阶段仍可能花费二次时间。输入含大量重复值、继续采用二路
Lomuto 划分时，互异元素下的期望证明也不适用。

\hypertarget{ux7406ux8bbaux4e0aux600eux6837ux9009ux62e9-k}{%
\subsection{\texorpdfstring{7. 理论上怎样选择
\(k\)}{7. 理论上怎样选择 k}}\label{ux7406ux8bbaux4e0aux600eux6837ux9009ux62e9-k}}

若目标是保持最优的渐近量级，取一个固定常数 \(k\) 即可：

\[
nk=O(n),\qquad n\lg(n/k)=\Theta(n\log n).
\]

取 \(k=O(\log n)\)，总界也保持在
\(O(n\log n)\)，与固定常数的渐近阶相同；取 \(k=n\) 时，上界则为
\(O(n^2)\)。

若进一步比较两阶段的成本常数，可以使用模型：

\[
F(k)=a\,n\ln(n/k)+b\,nk,\qquad a,b>0.
\]

把 \(k\) 暂当作连续变量，则：

\[
F'(k)=-\frac{an}{k}+bn=0
\quad\Longrightarrow\quad
k=\frac ab.
\]

这个模型给出的理想阈值为常数规模，具体取决于两阶段的成本系数。它不能确定一个适用于所有机器和程序的整数；渐近大
\(O\) 隐藏了这些系数，实际最优阈值仍需测量。

\hypertarget{ux5b9eux8df5ux4e2dux600eux6837ux9009ux62e9-k}{%
\subsection{\texorpdfstring{8. 实践中怎样选择
\(k\)}{8. 实践中怎样选择 k}}\label{ux5b9eux8df5ux4e2dux600eux6837ux9009ux62e9-k}}

固定实现与硬件，在同一批数据上比较几个小阈值。测试应覆盖随机、近乎有序、逆序及重复值较多的输入，并通过重复运行观察稳定耗时，减少随机枢轴波动对判断的干扰。

短数组上，插入排序的循环紧凑，局部访问连续；快速排序则要承担枢轴选择、划分和递归管理的开销。阈值决定何时切换两种算法，合适的取值还受元素大小、比较函数成本和缓存行为影响，需要通过测量确定。

\hypertarget{ux8bc1ux660eux4e2dux9700ux8981ux4fddux7559ux7684ux6761ux4ef6}{%
\subsection{9.
证明中需要保留的条件}\label{ux8bc1ux660eux4e2dux9700ux8981ux4fddux7559ux7684ux6761ux4ef6}}

\begin{itemize}
\tightlist
\item
  截止条件是「少于 \(k\)」，长度等于 \(k\) 时仍划分。
\item
  插入排序执行于整个数组；它的代价由全局剩余逆序对数控制。
\item
  小块内部无序，但块间顺序已经确定，否则 \(I=O(nk)\) 无法成立。
\item
  期望划分代价需要概率论证；理想均分的递归树只能提供直观解释。
\item
  小区间划分使用 \(O(1)\) 辅助空间，不等于整段递归算法也只用 \(O(1)\)
  空间。
\end{itemize}

\hypertarget{ux5341ux56dbux590dux4e60ux9898}{%
\section{十四、复习题}\label{ux5341ux56dbux590dux4e60ux9898}}

\begin{enumerate}
\def\labelenumi{\arabic{enumi}.}
\tightlist
\item
  多数元素的 \texttt{count}
  为什么不是实际出现次数？去掉存在性保证后应补哪一步？
\item
  Lomuto
  划分的四个区域是什么？为什么最后交换后还不能说左右两侧已经有序？
\item
  为什么互异元素下，枢轴秩的均匀分布足以分析任意固定输入？
\item
  快速排序的期望递推为什么能通过相邻规模作差消掉求和？
\item
  递归层数为对数，为什么 Quickselect 的时间仍为线性？
\item
  自底向上建堆为什么是线性时间？查看堆顶与取出堆顶的成本有何区别？
\item
  \(O(n+K)\) 为什么未必是关于二进制编码长度的多项式？
\item
  比较排序的叶子为什么至少有 \(n!\)
  个？下界为什么允许某些输入在线性时间内完成？
\item
  最大子数组的两个状态各自覆盖哪些候选区间？
\item
  多路归并为什么只需为每张表保留一个候选元素？
\item
  混合排序第一阶段之后，为什么逆序对只能留在小块内部？
\item
  怎样在允许随机坏划分的情况下，证明第一阶段的期望代价？
\end{enumerate}

\hypertarget{ux9644-ux8d44ux6599ux4e0bux8f7d}{%
\section{附 · 资料下载}\label{ux9644-ux8d44ux6599ux4e0bux8f7d}}

主要依据为课程课件《04-advanced\_sorting》、本次课堂转写《算法课排序复杂度与决策树证明》，以及当天课堂摘要。三道作业的编号、输入要求和目标复杂度均以课件末尾题目为准。

\begin{itemize}
\tightlist
\item
  \href{https://redmoon-2333.github.io/assets/attach/algorithms-day4/day4_algorithms.py}{课堂算法与三道作业的
  Python 实现}
\item
  \href{https://redmoon-2333.github.io/assets/attach/algorithms-day4/verify_day4.py}{算法与边界条件验证脚本}
\item
  \href{https://redmoon-2333.github.io/assets/attach/algorithms-day4/make_figures.py}{理论比较次数图的生成脚本}
\end{itemize}

单页手写提纲用于快速回顾课堂主线，详细推导和作业解答见本文：

单页手写速览另存为《Day4-单页手写速览.pdf》，可与本册配合阅读。

\end{document}
